%
% parabel.tex -- Exhaustion der Parabel
%
% (c) 2021 Prof Dr Andreas Müller, OST Ostschweizer Fachhochschule
%
\documentclass[tikz]{standalone}
\usepackage{amsmath}
\usepackage{times}
\usepackage{txfonts}
\usepackage{pgfplots}
\usepackage{csvsimple}
\usetikzlibrary{arrows,intersections,math}
\definecolor{darkred}{rgb}{0.8,0,0}
\begin{document}
\def\skala{6}
\def\n{5}
\begin{tikzpicture}[>=latex,thick,scale=\skala]

\foreach \i in {1,...,\n}{
	\pgfmathparse{\i/\n}
	\xdef\x{\pgfmathresult}
	\fill[color=blue!10]
		({\x-1/\n},{(\x-1/\n)*(\x-1/\n)})
		--
		({\x},{\x*\x})
		--
		({\x},1)
		--
		++({-1/\n},0)
		-- cycle;
	\draw[color=blue,line width=0.4pt]
		({\x-1/\n},{(\x-1/\n)*(\x-1/\n)})
		--
		({\x},{\x*\x});
	\node[color=blue] at
		({\x-0.5/\n},{0.5+0.25*((\x-1/\n)*(\x-1/\n)+\x*\x)}) {$F_\i$};
}

\foreach \i in {0,...,\n}{
	\pgfmathparse{\i/\n}
	\xdef\x{\pgfmathresult}
	\draw[color=blue,line width=0.4pt] ({\x},{\x*\x}) -- (\x,1);
	\draw ({\i/\n},{-0.05/\skala}) -- ++(0,{0.1/\skala});
}
\draw[color=blue,line width=0.4pt] (0,1) -- (1,1);

\pgfmathparse{\n-1}
\xdef\N{\pgfmathresult}
\foreach \i in {1,...,\N}{
	\node at ({\i/\n},{-0.05/\skala}) [below] {$\frac{\i}{\n}$};
}
\node at (0,{-0.05/\skala}) [below] {$0$};
\node at (1,{-0.05/\skala}) [below] {$1$};


\draw[color=darkred,line width=1.2pt] plot[domain=0:1,smooth] ({\x},{\x*\x});

\foreach \i in {0,...,\n}{
	\pgfmathparse{\i/\n}
	\xdef\x{\pgfmathresult}
	\fill[color=blue] ({\x},{\x*\x}) circle[radius={0.05/\skala}];
}

\draw[->] (-0.01,0) -- (1.05,0) coordinate[label={$x$}];
\draw[->] (0,-0.01) -- (0,1.05) coordinate[label={$y$}];

\end{tikzpicture}
\end{document}

